\chapter{Innovación y modelo de negocio}\label{appendix:innovacion}

Este apéndice complementa el desarrollo técnico descrito en el Capítulo~\ref{chapter:desarrollo} con
una perspectiva de innovación y viabilidad empresarial. Se analiza el mercado de la documentación
clínica asistida por \englishText{IA}, se define el modelo de negocio mediante el lienzo de modelo de
negocio (\englishText{Business Model Canvas}) y se describen el producto mínimo viable, el
\textit{roadmap} de evolución y los aspectos legales relevantes para una eventual explotación
comercial del sistema.

\section{Estudio del Mercado}\label{section:innovacion_mercado}

El mercado de la transcripción y documentación clínica asistida por voz se encuentra en expansión
acelerada a nivel internacional, impulsado por la sobrecarga administrativa del personal sanitario y
la madurez reciente de los modelos \englishText{ASR} y \englishText{LLM} descrita en el
Capítulo~\ref{chapter:marco_teorico_y_referencial}. No obstante, el segmento hispanohablante presenta
una oportunidad relativamente desatendida: la mayoría de soluciones comerciales consolidadas
están optimizadas para inglés y para los
flujos de trabajo y terminología del sistema sanitario de habla inglesa, con una adaptación limitada al
español médico, a \englishText{SNOMED CT} en su edición nacional y a \englishText{CIE-10-ES}.

\subsection{Segmentos de Clientes}\label{section:innovacion_segmentos}

El producto se orienta a un enfoque mixto \textit{B2B} y \textit{B2C}:

\begin{itemize}
    \item \textbf{\englishText{B2B} -- centros sanitarios:} hospitales, clínicas y centros médicos
    (públicos y privados) que buscan reducir la carga administrativa, optimizar flujos de trabajo y
    estandarizar la documentación clínica a escala de servicio o departamento, integrándose con sus
    sistemas de Historia Clínica Electrónica (\englishText{HCE}).
    \item \textbf{\englishText{B2C} -- profesionales independientes:} dueños de consultas privadas
    y pequeñas consultas que necesitan una herramienta ágil, de coste predecible y sin integración
    compleja, para generar informes normalizados sin depender de personal administrativo.
\end{itemize}

\subsection{Competidores}\label{section:innovacion_competidores}

El panorama competitivo puede agruparse en tres categorías:

\begin{itemize}
    \item \textbf{Asistentes de documentación clínicos basados en \englishText{IA} (mercado anglosajón):}
    soluciones como \englishText{Nuance Dragon Medical One / DAX}, \englishText{Abridge},
    \englishText{Suki AI} o \englishText{Corti} dominan el mercado de habla inglesa. Ofrecen
    transcripción ambiental y generación de notas \englishText{SOAP}, pero su despliegue depende de
    infraestructura en la nube de terceros y su adaptación al español clínico y a las ontologías
    locales (\englishText{SNOMED CT} edición española, \englishText{CIE-10-ES}) es limitada o
    inexistente.
    \item \textbf{Sistemas de dictado genéricos:} herramientas de transcripción de propósito general
    (por ejemplo, las integradas en suites ofimáticas) que producen texto libre sin estructuración
    clínica, normalización terminológica ni exportación interoperable, trasladando todo el trabajo de
    codificación al profesional.
    \item \textbf{Módulos nativos de \englishText{HCE}:} algunos proveedores de historia clínica
    electrónica incorporan dictado por voz como funcionalidad adicional, pero típicamente sin
    extracción semántica de entidades ni mapeo automático a estándares de codificación.
\end{itemize}

La diferenciación del sistema propuesto frente a estos competidores reside en tres elementos descritos
en el Capítulo~\ref{chapter:desarrollo}: 
\begin{enumerate}
\item Una capa intermedia de \englishText{NER} clínico en
español que normaliza explícitamente a \englishText{SNOMED CT} y \englishText{CIE-10-ES} antes de la
generación del informe, reduciendo el riesgo de alucinación frente a sistemas que delegan toda la
estructuración en el \englishText{LLM}.
\item La posibilidad de despliegue \textit{on-premise} con
modelos abiertos (Ollama, \texttt{llama3.1:8b}), evitando el envío de datos clínicos sensibles a
\englishText{APIs} de terceros.
\item La salida nativa en formato \englishText{HL7 FHIR}, lista para
integración con sistemas de \englishText{HCE} del Sistema Nacional de Salud.
\end{enumerate}

\section{Propuesta de Valor}\label{section:innovacion_propuesta_valor}

Facilidad de la generación de informes clínicos a partir de dictados médicos, reduciendo el tiempo
dedicado a documentación manual y minimizando errores de transcripción, con un informe final
normalizado contra estándares clínicos reconocidos (\englishText{SNOMED CT}, \englishText{CIE-10-ES})
y exportado en formato interoperable \englishText{HL7 FHIR}, lo que mejora su integración con
sistemas de historia clínica electrónica. La opción de despliegue local de los componentes de
\englishText{IA} (sin envío de audio o texto clínico a servicios externos de pago por uso) constituye
además un argumento diferencial frente a soluciones puramente \textit{cloud} en materia de privacidad
y soberanía del dato.

\section{Recursos Clave}\label{section:innovacion_recursos}

\begin{itemize}
    \item La aplicación cliente (captura de audio) y la pasarela de orquestación del \textit{pipeline}
    (Sección~\ref{section:arquitectura_sistema}).
    \item Los modelos de \englishText{ASR}, \englishText{NER} y el modelo de lenguaje (\englishText{LLM})
    encargado de estructurar y redactar el informe clínico.
    \item Un servicio de terminologías (\englishText{Snowstorm}) para mapear conceptos a estándares
    médicos (\englishText{SNOMED CT} y \englishText{CIE-10-ES}).
    \item El conjunto de evaluación y la metodología de medición (\englishText{WER}, F1 de
    \englishText{NER}, precisión de normalización, validación \englishText{FHIR} y juicio
    \englishText{LLM-as-judge}) desarrollados en el Capítulo~\ref{chapter:evaluacion}, que constituyen
    una ventaja competitiva para auditar y mejorar la calidad clínica del producto frente a
    competidores que no publican métricas de fiabilidad.
    \item Equipo de desarrollo y evaluación de la aplicación.
    \item Servicio de hosting donde vivirá el backend de la aplicación.
\end{itemize}

\section{Actividades Clave}\label{section:innovacion_actividades}

\begin{itemize}
    \item Mantenimiento y mejora continua del \textit{pipeline} de \englishText{IA} (re-entrenamiento
    o sustitución de modelos \englishText{ASR}/\englishText{NER}/\englishText{LLM} conforme evolucione
    el estado del arte, Sección~\ref{section:trabajos_futuros}).
    \item Validación clínica continua mediante revisión por profesionales sanitarios o agentes de IA y actualización
    del \textit{gold standard} de evaluación.
    \item Integración técnica con sistemas de \englishText{HCE} de clientes \textit{B2B} (adaptación de
    perfiles \englishText{FHIR}, mapeos específicos por centro).
    \item Cumplimiento normativo continuo (protección de datos clínicos, certificación como producto
    sanitario si aplica, véase Sección~\ref{section:innovacion_legal}).
    \item Soporte técnico y comercial diferenciado por segmento de cliente.
\end{itemize}

\section{Socios y Alianzas Clave}\label{section:innovacion_alianzas}

Como aliados clave se incluyen proveedores de \textit{hosting} e infraestructura, necesarios para
garantizar la disponibilidad y operación del sistema. En una fase de expansión, podrían incorporarse
proveedores de servicios de
transcripción especializados como \englishText{Google Chirp} o \englishText{Amazon Transcribe Medical}
para una transcripción inicial de mayor calidad en escenarios donde el cliente acepte el procesamiento
en la nube. Asimismo, resultan estratégicas las alianzas con centros sanitarios piloto para la
validación clínica del sistema y con el \englishText{Barcelona Supercomputing Center (BSC)}, cuyos
modelos y corpus en español (Sección~\ref{section:innovacion_mercado}) sustentan el componente de
\englishText{NER} clínico del sistema.

\section{Estructura de Costes Prevista}\label{section:innovacion_costes}

La estructura de costos se compone principalmente de:

\begin{itemize}
    \item \textbf{Adquisición de clientes:} inversión en marketing y equipo comercial, especialmente
    relevante en las etapas iniciales del proyecto.
    \item \textbf{Infraestructura tecnológica:} servidores, bases de datos, almacenamiento y cómputo
    (\englishText{GPU}) para inferencia de los modelos \englishText{ASR}/\englishText{NER}/\englishText{LLM}
    y para la instancia de \englishText{Snowstorm}, conforme a los requisitos descritos en
    Sección~\ref{section:gestion_herramientas}.
    \item \textbf{Mantenimiento y escalabilidad:} evolución de la plataforma, soporte técnico y
    actualización de modelos.
    \item \textbf{Cumplimiento normativo:} auditorías de protección de datos y, en su caso, costes de
    certificación regulatoria (Sección~\ref{section:innovacion_legal}).
\end{itemize}

En el despliegue \textit{on-premise}, los costes de inferencia se concentran en hardware local
(\englishText{CPU}/\englishText{GPU}) en lugar de en consumo de \englishText{APIs} de pago por uso, lo
que reduce el coste marginal por informe generado a medida que crece el volumen de uso, a costa de una
mayor inversión inicial en infraestructura.

\section{Vías de Ingresos}\label{section:innovacion_ingresos}

El modelo de ingresos se divide en \textit{B2B} y \textit{B2C}. En \textit{B2B}, se ofrecen contratos
de mantenimiento y soporte, ya que los informes pueden requerir adaptaciones según el centro médico
(perfiles \englishText{FHIR} específicos, integración con su \englishText{HCE}, diferentes formatos de salida).
En \textit{B2C}, se
propone una suscripción mensual o anual con acceso a la plataforma y generación de informes médicos
normalizados y estructurados. Como vías de financiación para las primeras fases, se contemplan rondas
de capital semilla orientadas a \textit{healthtech}, líneas de ayuda pública a la innovación digital en
salud, y programas piloto remunerados con centros sanitarios colaboradores que permitan financiar la
validación clínica a la vez que se recopila evidencia de uso real.

\section{Canales de Distribución}\label{section:innovacion_canales}

Los canales de distribución serán similares para ambos segmentos, al tratarse del mismo mercado
objetivo. Se contempla el uso de una página web corporativa, redes sociales, campañas de marketing
digital, así como la participación en ferias y eventos de tecnología y salud.

\section{Relación con los Clientes}\label{section:innovacion_relacion_clientes}

Para el segmento \textit{B2B}, se plantea una relación basada en un equipo de ventas y soporte técnico,
encargado de la captación, seguimiento y mantenimiento de los contratos con centros médicos. En el
segmento \textit{B2C}, la relación se gestiona mediante un canal de soporte centralizado, donde los
usuarios pueden reportar incidencias y recibir asistencia técnica.

\section{Descripción del Producto Mínimo Viable}\label{section:innovacion_mvp}

El sistema desarrollado y evaluado en este trabajo (Capítulos~\ref{chapter:desarrollo}
y~\ref{chapter:evaluacion}) constituye el producto mínimo viable: una aplicación cliente de captura de
audio, una pasarela de orquestación, y un \textit{pipeline} completo de transcripción \englishText{ASR},
extracción de entidades clínicas (\englishText{NER}), normalización a \englishText{SNOMED CT}/
\englishText{CIE-10-ES}, generación de informe mediante \englishText{LLM} y exportación en
\englishText{HL7 FHIR}. Como se concluye en el Capítulo~\ref{chapter:conclusiones}, el \englishText{MVP}
demuestra viabilidad técnica de extremo a extremo, pero su nivel de fiabilidad clínica actual
($1{,}55$--$2{,}90$ sobre 5 en evaluación \englishText{end-to-end}) requiere que el informe generado se
emplee como borrador asistido sujeto a revisión y firma de un profesional sanitario, y no como
documento clínico definitivo sin supervisión.

\section{Roadmap de Desarrollo}\label{section:innovacion_roadmap}

A partir de las líneas de trabajo futuro identificadas en la
Sección~\ref{section:trabajos_futuros}, se propone el siguiente \textit{roadmap} orientado a producto:

\begin{itemize}
    \item \textbf{Corto plazo:} mitigación de la propagación de errores entre etapas del
    \textit{pipeline} mediante verificación intermedia, y mejora del módulo \englishText{NER} (expansión
    de siglas y abreviaturas clínicas, reducción de la sobre-segmentación de \textit{spans}), como
    requisito previo a cualquier piloto con centros sanitarios.
    \item \textbf{Medio plazo:} normalización terminológica sensible al contexto mediante enlace de
    entidades (\textit{entity linking}), reducción de alucinaciones del informe \englishText{LLM}
    mediante técnicas de generación aumentada por recuperación (\englishText{RAG}) o verificación
    post-generación, e incorporación de integraciones \englishText{FHIR} específicas con \englishText{HCE}
    de los primeros clientes piloto.
    \item \textbf{Largo plazo:} evaluación y despliegue sobre infraestructura de mayor capacidad que
    permita emplear modelos \englishText{LLM} de mayor escala, ampliación a otras especialidades
    clínicas más allá del conjunto evaluado.
\end{itemize}

% \section{Aspectos Legales}\label{section:innovacion_legal}

% Al tratarse de un sistema que procesa datos de salud, su explotación comercial debe abordar, como
% mínimo, los siguientes aspectos:

% \begin{itemize}
%     \item \textbf{Protección de datos:} cumplimiento del Reglamento General de Protección de Datos
%     (\englishText{RGPD}) y de la Ley Orgánica 3/2018 (\englishText{LOPDGDD}) en el tratamiento de datos
%     de salud, considerados categoría especial. El despliegue \textit{on-premise} descrito en
%     Sección~\ref{section:innovacion_propuesta_valor} facilita este cumplimiento al evitar la
%     transferencia de datos clínicos a terceros, aunque dicha transferencia sí debe regularse mediante
%     los acuerdos correspondientes en los escenarios donde se empleen servicios de transcripción en la
%     nube (Sección~\ref{section:innovacion_alianzas}).
%     \item \textbf{Reglamento de Inteligencia Artificial de la Unión Europea:} un sistema que asiste en
%     la documentación clínica puede quedar comprendido en el ámbito de aplicaciones de alto riesgo
%     vinculadas a la salud, lo que conllevaría obligaciones de gestión de riesgos, calidad de los datos
%     de entrenamiento, trazabilidad y supervisión humana, coherentes con el uso del informe generado
%     como borrador asistido descrito en la Sección~\ref{section:innovacion_mvp}.
%     \item \textbf{Producto sanitario (\englishText{MDR}):} en función del uso final, una herramienta que
%     estructura y codifica información clínica de forma automatizada podría requerir evaluación de su
%     clasificación bajo el Reglamento (UE) 2017/745 de productos sanitarios, especialmente si su salida
%     se integra directamente en la \englishText{HCE} sin revisión humana intermedia.
%     \item \textbf{Propiedad intelectual y licencias:} los modelos de terceros empleados (modelos
%     \englishText{ASR}, \englishText{NER} del \englishText{BSC}, modelos \englishText{LLM} servidos vía
%     Ollama) deben revisarse en cuanto a sus licencias de uso comercial, así como las condiciones de uso
%     de \englishText{Snowstorm} y de las ontologías \englishText{SNOMED CT} y \englishText{CIE-10-ES}.
% \end{itemize}
